% =====================================================================
% CODE LISTINGS FOR THE IMPLEMENTATION SECTIONS OF SCENARIOS 1 - 4
%
% HOW TO USE THIS FILE
%   1. Copy the PREAMBLE BLOCK below into your document preamble (once).
%   2. Paste each scenario's block into that scenario's
%      "Scenario Implementation" subsection, after the existing paragraph.
%
% Every listing here is taken verbatim from the working code, with only
% blank lines and repeated boilerplate removed to keep the report readable.
% =====================================================================


% =====================================================================
% PREAMBLE BLOCK - paste this ONCE, in your document preamble
% =====================================================================
%
% \usepackage{listings}
% \usepackage{xcolor}
%
% \definecolor{codebg}{rgb}{0.97,0.97,0.97}
% \definecolor{codecomment}{rgb}{0.00,0.45,0.00}
% \definecolor{codekeyword}{rgb}{0.10,0.10,0.65}
% \definecolor{codestring}{rgb}{0.60,0.10,0.40}
% \definecolor{coderule}{rgb}{0.65,0.65,0.65}
%
% % listings has no built-in YAML support, so we define a minimal one.
% \lstdefinelanguage{yaml}{
%   keywords={apiVersion,kind,metadata,spec,template,selector,matchLabels,
%             labels,replicas,containers,image,name,command,env,value,ports,
%             nodes,role,nodeSelector,affinity,podAntiAffinity,tolerations,
%             readinessProbe,livenessProbe,resources,services,networks,
%             build,depends_on,driver,true,false},
%   sensitive=true,
%   comment=[l]{\#},
%   morestring=[b]",
%   morestring=[b]'
% }
%
% \lstdefinestyle{scenariocode}{
%   backgroundcolor=\color{codebg},
%   basicstyle=\ttfamily\footnotesize,
%   commentstyle=\color{codecomment}\itshape,
%   keywordstyle=\color{codekeyword}\bfseries,
%   stringstyle=\color{codestring},
%   numbers=left,
%   numberstyle=\tiny\color{coderule},
%   numbersep=8pt,
%   frame=single,
%   rulecolor=\color{coderule},
%   breaklines=true,
%   breakatwhitespace=false,
%   captionpos=b,
%   keepspaces=true,
%   showspaces=false,
%   showstringspaces=false,
%   showtabs=false,
%   tabsize=4,
%   xleftmargin=16pt,
%   framexleftmargin=16pt
% }
% \lstset{style=scenariocode}
%
% =====================================================================



% #####################################################################
% SCENARIO 1 - THE SILENT MEMORY LEAK
% Paste into the "Scenario Implementation" subsection of Scenario 1.
% #####################################################################

All four scenarios share the same telemetry pipeline. A single
\texttt{TracerProvider} is created, given a service name, and attached to a
\texttt{BatchSpanProcessor} that exports finished spans to Jaeger over OTLP.
Listing~\ref{lst:otel-setup} shows this shared setup.

\begin{lstlisting}[language=Python, caption={Shared OpenTelemetry pipeline setup (\texttt{scenario\_1.py}).}, label={lst:otel-setup}]
from opentelemetry import trace
from opentelemetry.sdk.trace import TracerProvider
from opentelemetry.sdk.trace.export import BatchSpanProcessor
from opentelemetry.exporter.otlp.proto.grpc.trace_exporter import OTLPSpanExporter
from opentelemetry.sdk.resources import Resource
from opentelemetry.trace.status import Status, StatusCode

# 1. Setup the Pipeline
resource = Resource(attributes={"service.name": "multi-level-app"})
provider = TracerProvider(resource=resource)
processor = BatchSpanProcessor(
    OTLPSpanExporter(endpoint="http://localhost:4317", insecure=True))
provider.add_span_processor(processor)
trace.set_tracer_provider(provider)
tracer = trace.get_tracer(__name__)
\end{lstlisting}

The failure itself is a hard threshold inside the Level 2 database span. The
simulated processing time grows on every iteration, and once it passes
$0.8$ seconds the span is marked with an \texttt{OutOfMemoryError} and returns a
failure to Level 1, which in turn marks its own span as failed. This is what
produces the paired red spans in the exported trace data.

\begin{lstlisting}[language=Python, caption={The memory-leak threshold and the cascade into Level 1 (\texttt{scenario\_1.py}).}, label={lst:s1-threshold}]
# --- LEVEL 2: The Simulated Database (Memory Leak Degradation) ---
def level_2_database_call(current_latency):
    with tracer.start_as_current_span("level-2-db-query") as span:
        span.set_attribute("db.type", "sql")
        time.sleep(current_latency)          # degrading processing time

        # Threshold: past 0.8 s the server runs out of RAM and crashes
        if current_latency > 0.8:
            error_message = "FATAL: OutOfMemoryError - System Degraded"
            span.add_event("exception", {"exception.message": error_message})
            span.set_status(Status(StatusCode.ERROR, error_message))
            return False                     # Failed
        return True                          # Healthy

# --- LEVEL 1: The Simulated Frontend API ---
def level_1_api_request(current_latency):
    with tracer.start_as_current_span("level-1-api-endpoint") as span:
        span.set_attribute("http.method", "GET")
        success = level_2_database_call(current_latency)
        if not success:
            span.set_status(Status(StatusCode.ERROR,
                "Request timed out due to downstream degradation"))
\end{lstlisting}

The degradation is driven by a $60$-step loop. Latency begins at $100$
milliseconds and grows by $15$ milliseconds per step, with a small random
jitter added so the series resembles real measurements rather than a perfectly
straight line.

\begin{lstlisting}[language=Python, caption={The time-series degradation loop (\texttt{scenario\_1.py}).}, label={lst:s1-loop}]
base_latency = 0.1        # Starts fast and healthy (100 ms)
degradation_step = 0.015  # Latency increases by 15 ms every step

for i in range(60):
    base_current = base_latency + (i * degradation_step)
    jitter = random.uniform(-0.02, 0.02)     # real-world fluctuation
    actual_latency = max(0.01, base_current + jitter)
    level_1_api_request(actual_latency)

provider.force_flush()
\end{lstlisting}



% #####################################################################
% SCENARIO 2 - CASCADING TIMEOUT (TRAFFIC SPIKE)
% Paste into the "Scenario Implementation" subsection of Scenario 2.
% #####################################################################

The defining feature of this scenario is that the two levels disagree about
whether anything went wrong. Listing~\ref{lst:s2-timeout} shows why. Level 2
never raises an error: it simply adds the queue delay to its own fast
processing time and reports success. Level 1 applies a strict
\texttt{if total\_db\_time > 2.0} rule and abandons the request, marking only
its own span as failed. The result is a trace whose backend is green and whose
frontend is red, which is exactly the symptom-versus-cause distinction the
Explainable AI must resolve.

\begin{lstlisting}[language=Python, caption={A healthy but slow database, and the strict frontend timeout that fails above it (\texttt{scenario\_2.py}).}, label={lst:s2-timeout}]
# --- LEVEL 2: The Database (Getting overloaded by the queue) ---
def level_2_database_call(queue_delay):
    with tracer.start_as_current_span("level-2-db-query") as span:
        span.set_attribute("db.type", "sql")
        # Processing is fast (0.1 s), but the line to get in is huge
        actual_processing_time = 0.1
        total_time = actual_processing_time + queue_delay
        time.sleep(total_time)
        # No error here. The database is healthy, just backed up.
        return total_time

# --- LEVEL 1: The Frontend API (strict 2.0 second timeout) ---
def level_1_api_request(queue_delay):
    with tracer.start_as_current_span("level-1-api-endpoint") as span:
        span.set_attribute("http.method", "POST")
        span.set_attribute("http.route", "/checkout")
        total_db_time = level_2_database_call(queue_delay)

        # --- THE CASCADING TIMEOUT ---
        if total_db_time > 2.0:
            error_msg = ("HTTP 504: Gateway Timeout - "
                         "Downstream database took too long")
            span.add_event("exception", {"exception.message": error_msg})
            span.set_status(Status(StatusCode.ERROR, error_msg))
            return False
        return True
\end{lstlisting}

The traffic spike is injected at the midpoint of the run. For the first $30$
iterations the queue is empty and every checkout succeeds. From iteration $30$
onwards the queue delay grows by $0.2$ seconds per minute, because requests
arrive faster than the database can clear them. The frontend therefore begins
failing once the accumulated delay pushes the total past its two-second budget.

\begin{lstlisting}[language=Python, caption={Injecting the traffic spike at minute 30 (\texttt{scenario\_2.py}).}, label={lst:s2-spike}]
queue_delay = 0.0

for i in range(60):                # Simulate 60 minutes
    if i == 30:
        print("MASSIVE TRAFFIC SPIKE! 5,000 users clicked 'Buy' at once!")
    if i >= 30:
        # Requests arrive faster than they finish, so the queue grows
        queue_delay += 0.2

    jitter = random.uniform(-0.02, 0.02)
    current_delay = max(0.0, queue_delay + jitter)
    success = level_1_api_request(current_delay)
\end{lstlisting}



% #####################################################################
% SCENARIO 3 - NETWORK CHAOS (PACKET LOSS AND CONGESTION)
% Paste into the "Scenario Implementation" subsection of Scenario 3.
% #####################################################################

Unlike the first two scenarios, this one requires real machines talking over a
real network, so the services were containerised and placed on a private Docker
network. Listing~\ref{lst:s3-compose} shows the three containers: the Jaeger
trace store, the Level 2 backend and the Level 1 frontend. Note that this
scenario exports traces to \texttt{jaeger:4317}, the container hostname, rather
than to \texttt{localhost}.

\begin{lstlisting}[language=yaml, caption={The private Docker network and its three containers (\texttt{docker-compose.yml}).}, label={lst:s3-compose}]
services:
  jaeger:                                  # 1. The Trace Database
    image: jaegertracing/all-in-one:latest
    ports:
      - "16686:16686"
      - "4317:4317"
    networks:
      - ainetwork

  level-2:                                 # 2. The Backend Service
    build: .
    command: tail -f /dev/null             # keep the container awake
    networks:
      - ainetwork

  level-1:                                 # 3. The Frontend API Service
    build: .
    command: tail -f /dev/null
    networks:
      - ainetwork
    depends_on:
      - level-2
      - jaeger

networks:                                  # a private virtual network
  ainetwork:
    driver: bridge
\end{lstlisting}

The application code makes genuine HTTP calls across that network with a strict
$0.5$ second budget. When congestion delays or drops the packets, the request
raises a \texttt{RequestException}, which is captured on the span as an error.
Nothing inside the application is broken; the wire between the services is.

\begin{lstlisting}[language=Python, caption={Real HTTP requests with a strict timeout, recording network failures as error spans (\texttt{network\_test.py}).}, label={lst:s3-request}]
def make_network_request(attempt_number):
    with tracer.start_as_current_span("http-request-to-level-2") as span:
        span.set_attribute("attempt.number", attempt_number)
        try:
            # Call the Level 2 container over the Docker network
            start_time = time.time()
            response = requests.get("http://level-2:8000", timeout=0.5)
            duration = time.time() - start_time
            span.set_attribute("http.status_code", response.status_code)

        except requests.exceptions.RequestException as e:
            # If packet loss is too high, the request fails
            error_msg = f"Network Failure: {str(e)}"
            span.add_event("exception", {"exception.message": error_msg})
            span.set_status(Status(StatusCode.ERROR, error_msg))
\end{lstlisting}

The congestion itself is produced by \texttt{iperf3}, which floods the private
network with $100$ Mbit/s of UDP traffic for $30$ seconds. The flood is started
first and the request loop immediately afterwards, so that the two overlap.

\begin{lstlisting}[language=bash, caption={Generating the network flood and running the experiment underneath it.}, label={lst:s3-flood}]
# Terminal 1 - Level 2: the web server under test
docker compose exec level-2 python -m http.server 8000

# Terminal 2 - Level 2: the iperf server that receives the flood
docker compose exec level-2 iperf3 -s

# Terminal 3 - Level 1: saturate the network with 100 Mbit/s UDP for 30 s
docker compose exec level-1 iperf3 -c level-2 -u -b 100M -t 30

# Terminal 4 - Level 1: fire the traced requests into the congestion
docker compose exec level-1 python3 network_test.py
\end{lstlisting}



% #####################################################################
% SCENARIO 4 - KUBERNETES FAULT TOLERANCE AND REROUTING
% Paste into the "Scenario Implementation" subsection of Scenario 4.
% #####################################################################

The cluster is declared as five machines, as shown in
Listing~\ref{lst:s4-cluster}. Labels separate the machine that carries the
observability tools from the three machines that run the order backend, which
guarantees that the fault we inject can never destroy the instrument measuring
it.

\begin{lstlisting}[language=yaml, caption={The five-machine cluster definition (\texttt{k8s/kind-cluster.yaml}).}, label={lst:s4-cluster}]
kind: Cluster
apiVersion: kind.x-k8s.io/v1alpha4
name: rca4

nodes:
  - role: control-plane          # etcd, apiserver, scheduler, controllers
    extraPortMappings:
      - containerPort: 30080     # api-gateway, published to the host
        hostPort: 30080
      - containerPort: 30686     # Jaeger UI and query API
        hostPort: 30686

  - role: worker                 # observability + entry point; never killed
    labels:
      rca4-role: infra

  - role: worker                 # the application tier: one of these is shot
    labels:
      rca4-role: app
  - role: worker
    labels:
      rca4-role: app
  - role: worker
    labels:
      rca4-role: app
\end{lstlisting}

Listing~\ref{lst:s4-deploy} contains the three settings that make self-healing
observable inside a short experiment. The anti-affinity rule spreads the copies
one per machine; the tolerations shorten the wait before a copy on an
unreachable machine is replaced, from the Kubernetes default of $300$ seconds
down to $10$; and the environment variables fix the capacity of each copy and
the length of its warm-up.

\begin{lstlisting}[language=yaml, caption={Replica placement, fast failover and the capacity settings (\texttt{k8s/20-order-backend.yaml}).}, label={lst:s4-deploy}]
spec:
  replicas: 3
  template:
    spec:
      nodeSelector:
        rca4-role: app
      affinity:
        podAntiAffinity:                 # one copy per machine
          preferredDuringSchedulingIgnoredDuringExecution:
            - weight: 100
              podAffinityTerm:
                labelSelector:
                  matchLabels:
                    app: order-backend
                topologyKey: kubernetes.io/hostname
      tolerations:                       # default is 300 s; we want 10 s
        - key: node.kubernetes.io/unreachable
          operator: Exists
          effect: NoExecute
          tolerationSeconds: 10
      terminationGracePeriodSeconds: 0   # abrupt loss, not a polite shutdown
      containers:
        - name: order-backend
          env:
            - name: WORKER_SLOTS         # 2 slots / 100 ms = 20 requests/s
              value: "2"
            - name: SERVICE_TIME_MS
              value: "100"
            - name: WARMUP_SECONDS       # a new copy is useless this long
              value: "15"
\end{lstlisting}

The capacity limit is enforced in the application itself, as shown in
Listing~\ref{lst:s4-capacity}. A semaphore admits only \texttt{WORKER\_SLOTS}
requests at a time; everything else waits. The waiting time is measured and
written onto the span, which makes the growth of the queue directly visible in
the telemetry and is the measurement the brownout is read from.

\begin{lstlisting}[language=Python, caption={The bounded-concurrency limit that converts lost capacity into measurable queueing (\texttt{app/order\_backend.py}).}, label={lst:s4-capacity}]
# Capacity per replica = WORKER_SLOTS / SERVICE_TIME_MS = 20 requests/s
slots = threading.Semaphore(WORKER_SLOTS)

def _handle_order(self):
    # Continue the trace the API gateway started
    ctx = propagator.extract(
        carrier={k.lower(): v for k, v in self.headers.items()})

    with tracer.start_as_current_span("order-backend.process", context=ctx) as span:
        span.set_attribute("k8s.pod.name", POD_NAME)
        span.set_attribute("k8s.node.name", NODE_NAME)

        # Time spent waiting for a free slot IS the brownout signal
        queued_at = time.time()
        slots.acquire()
        queue_wait = time.time() - queued_at
        try:
            time.sleep(SERVICE_TIME_MS / 1000.0)
        finally:
            slots.release()

        span.set_attribute("backend.queue_wait_ms", round(queue_wait * 1000, 2))
\end{lstlisting}

The fault is injected by destroying the container that hosts one of the
machines. \texttt{docker kill} is used rather than a graceful stop so that the
machine disappears with no shutdown sequence at all, exactly like a pulled power
cable. The precise moment is recorded, because every later measurement and the
training label are expressed relative to it.

\begin{lstlisting}[language=Python, caption={Injecting the node failure and recording the exact moment (\texttt{scripts/inject\_chaos.py}).}, label={lst:s4-chaos}]
target_pod, target_node = pick_target(pods, args.target)

fault_at = time.time()
if args.mode == "node":
    # `docker kill` rather than `stop`: no SIGTERM, no graceful shutdown.
    # The machine simply ceases to exist as far as the cluster is concerned.
    run(["docker", "kill", target_node])
else:
    run(["kubectl", "-n", NAMESPACE, "delete", "pod", target_pod,
         "--grace-period=0", "--force"])

event = {"mode": args.mode, "target_pod": target_pod,
         "target_node": target_node, "fault_at": fault_at}
\end{lstlisting}

Roughly $12{,}000$ traces per run are then flattened into one row per second,
and each row is labelled with the number of seconds remaining until the system
is stable again. Listing~\ref{lst:s4-label} shows how that moment is located.
The rule compares against the measured baseline rather than against perfection,
and accepts $8$ healthy seconds out of $10$, because a healthy cluster still
produces the occasional slow request and the finish line must not depend on a
single unlucky packet.

\begin{lstlisting}[language=Python, caption={Locating the moment of restabilization, which produces the training label (\texttt{scripts/build\_dataset.py}).}, label={lst:s4-label}]
threshold = max(baseline_p95 * tolerance, baseline_p95 + 20)
err_threshold = max(base_err_rate * 2.0, 0.02)

def is_healthy(row):
    return (row["err_rate"] <= err_threshold
            and row["active_pods"] >= expected_replicas
            and row["p95_ms"] <= threshold)

# A strong majority, not every second: one noisy second must not move the
# finish line. But the window must *start* healthy, so the timestamp still
# marks a real transition.
needed = int(round(hold * hold_fraction))
for row in incident:
    start = row["t_rel"]
    window = [by_t[t] for t in range(start, start + hold) if t in by_t]
    if len(window) < hold or not is_healthy(window[0]):
        continue
    if sum(1 for w in window if is_healthy(w)) >= needed:
        return start, threshold, err_threshold

# Label for every second of the incident
row["seconds_to_restabilize"] = max(0, restabilized_at - row["t_rel"])
\end{lstlisting}

Finally the forecaster is trained and validated. Because one outage produces
only one recovery curve, the models are scored with \textit{leave-one-run-out}
testing: the model is trained on complete outages and asked to predict a
different outage it has never seen.

\begin{lstlisting}[language=Python, caption={Leave-one-run-out validation of the recovery forecaster (\texttt{scripts/forecast\_recovery.py}).}, label={lst:s4-train}]
FULL_FEATURES = ["p50_ms", "p95_ms", "max_ms", "mean_ms", "mean_queue_wait_ms",
                 "err_rate", "err_rate_roll5", "active_pods", "pods_missing",
                 "p95_ratio", "p95_slope_5s", "req_count", "time_since_fault"]
TARGET = "seconds_to_restabilize"

models = {
    "RandomForest": RandomForestRegressor(n_estimators=300, random_state=42),
    "GradientBoosting": GradientBoostingRegressor(random_state=42),
    "LinearRegression": LinearRegression(),
}

# Train on whole outages, test on an outage the model has never seen
splitter = GroupKFold(n_splits=min(len(runs), 5))
for name, model in models.items():
    maes = []
    for train_idx, test_idx in splitter.split(X, y, groups):
        model.fit(X[train_idx], y[train_idx])
        predicted = model.predict(X[test_idx])
        maes.append(mean_absolute_error(y[test_idx], predicted))
    print(f"{name}: MAE {np.mean(maes):.2f}s")
\end{lstlisting}

The same evaluation is then repeated with \texttt{time\_since\_fault} removed
from the feature list, forcing the model to rely only on the shape of the
latency and replica-count curves. Comparing the two results is what confirms the
model has learned the recovery dynamics rather than simply reading a clock.
